\documentclass[12pt]{article}

\usepackage{amsmath}
\usepackage[margin=1in]{geometry}
\usepackage{fancyhdr}
\usepackage{amsthm}
\newtheorem{lemma}{Lemma}
\usepackage{braket}
\usepackage{amssymb}
\usepackage{tikz}
\usepackage{bbm}
\usetikzlibrary{shapes.geometric}
\usepackage{enumerate}
\usepackage[shortlabels]{enumitem}


\pagestyle{fancy}
\fancyhead[l]{Hudson Benites}
\fancyhead[c]{Day 1}
\fancyfoot[c]{\thepage}
\renewcommand{\headrulewidth}{0.2pt}
\setlength{\headheight}{15pt}
\fancyhead[r]{\today}

\begin{document}
\begin{enumerate}[label=\Alph*.]
    \item
    \begin{enumerate}[label=\theenumi\Alph*]
        \item The smallest class of $f$ is $C^2$.
        \item $f(x)=\begin{cases}
            0 & \text{if } x \leq 0 \\
            x^{k+1} & \text{if } x > 0
        \end{cases}$.
        \item First, we show that each of the component functions of $f$ are themselves smooth. Let $h(x)=0$, which is constant. Therefore, $h^{(k)}(x)=0$, and as such is smooth. Let $g(x)=e^{-x^{-1}}$. It is an exponential, and therefore
              is smooth. We need only now consider $x=0$, where we show that $\lim_{x\rightarrow 0^-}f^{(k)}(x)=\lim_{x\rightarrow 0^+}f^{(k)}(x)$ for all $k\geq 0$. The $k$th derivative of $g(x)$ is \begin{equation}
              \frac{d^{(k)}}{dx^{(k)}}(-x^{-1})e^{-x^{-1}}.
              \end{equation} The $k$th derivative of $-x^{-1}$ is in the form $\pm \frac{1}{x^k}$, and therefore the $k$th derivative of $g(x)$ is $\pm \frac{e^{-x^{-1}}}{x^k}$. We assume without proof that $\lim_{x\rightarrow 0^+}\frac{e^{-x^{-1}}}{x^k}=0$. We have
              previously shown that $h^{(k)}(x)=0$, and $\lim_{x\rightarrow 0^-}{0}=0$. Therefore, we have shown $f$ to be smooth, and we are done.
        \item $f$ is not analytic.
            \begin{proof}
                We show that $f$ is not analytic by showing that its Taylor-Maclaurin series does not converge to $f$ on all of $\mathbb{R}$. The Taylor-Maclaurin series for  $e^{-x^{-1}}$ is
              $\sum_{n=0}^{\infty}\frac{(-1)^n}{n!x^n}$. When $x=0$, this sum does not exist, and as a result the Taylor-Maclaurin
              series of $f$ does not converge to $f$ on all of $\mathbb{R}$. Thus, we are done.
            \end{proof}
    \end{enumerate}
    
    \item
    \begin{enumerate}[label=\theenumi\Alph*]
        \item  For a positive-definite, symmetric, $m\times m$ matrix $G$, there exists a $\sqrt{G}$ such that $(\sqrt{G})^2 = G$.
            \begin{proof}
                Let $G$ be a symmetric, positive-definite, $m\times m$ matrix. By the spectral theorem, we know that
              $G$ has $m$ orthogonal (and thus linearly independent) eigenvectors, and a corresponding $m$ real eigenvalues.
              Furthermore, if $G$ has $m$ linearly independent eigenvectors, then the matrix $G$ diagonalizes as $PDP^{-1}$, where the columns
              of $P$ are the $m$ eigenvectors of $G$, and $D$ is a diagonal matrix whose entries are the corresponding eigenvalues of $G$. Being 
              positive-definite, $G$ has stricly positive eigenvalues, and therefore $D$ has contains only zeros and positive real numbers. As a result,
              we can construct a matrix $\sqrt{D}$, such that $({\sqrt{D}})^2=D$ as follows:
              \begin{equation*}
              \sqrt{D}=\begin{bmatrix}
              \sqrt{D_{1,1}} & 0 & \cdots & 0 \\
              0 & \sqrt{D_{2,2}} & \cdots & 0 \\
              \vdots & \vdots & \ddots & \vdots \\
              0 & 0 & \cdots & \sqrt{D_{m,m}}
              \end{bmatrix},
              \end{equation*}
              the entries of which are guaranteed to exist, as every entry of $D$ is greater than or equal to 0.
              With this in mind, consider the matrix $P\sqrt{D}P^{-1}$. First, we show that $P\sqrt{D}P^{-1}$ is symmetric by showing
              that it is equal to its tranpose. First note that because $G$ is symmetric, $PDP^{-1}=(P^{-1})^TDP^T$, and therefore
              $P^{-1}=P^{T}$. Taking the transpose of $P\sqrt{D}P^{-1}$ results in 
              \begin{equation}
                (P^{-1})^T\sqrt{D}P^T
              \end{equation}
              Using the fact that $P^{-1}=P{^T}$, (1) is equivalent to $P\sqrt{D}P^{-1}$, therefore $P\sqrt{D}P^{-1}$ is a symmetric matrix.
              Squaring this matrix, we see that it amounts to $P\sqrt{D}P^{-1}P\sqrt{D}P^{-1}$, which is indeed $G$. Thus, we are done.
            \end{proof}
        


        \item There exists an invertible linear transformation $f: \mathbb{R}\rightarrow \mathbb{R}$ such that $v\cdot_{G} w = f(v)\cdot f(w)$ for all $v,w \in \mathbb{R}$.
              \begin{proof}
                Let $f(v)=\sqrt{G}v$. First, we show that $f$ is injective, and assume that $f(v)=f(w)$ for some $v,w \in {\mathbb{R}}^m$. By the definition of $f$, this is equivalent to $\sqrt{G}v=\sqrt{G}w$.
                We now show that $\sqrt{G}$ has an inverse, by demonstrating that $(\sqrt{G})^{-1}=P(\sqrt{D})^{-1}P^{-1}$ from problem B.A. Note that since $D$ is a stricly positive diagonal matrix, $\sqrt{D}$ is as well, and is therefore invertible.
                Then, $P\sqrt{D}P^{-1}P(\sqrt{D})^{-1}P^{-1}$ is indeed $I$, and we have shown $(\sqrt{G})^{-1}$ to be $P(\sqrt{D})^{-1}P^{-1}$. Left multiplying by $(\sqrt{G})^{-1}$, we see that $v=w$, thus showing that $f$ is injective.\newline
                Next, we show that $f$ is surjective, and let $y$ be some element of $\mathbb{R}^m$. Consider $x=(\sqrt{G})^{-1}y$, which is also in $\mathbb{R}^m$.
                Then, $f(x)=\sqrt{G}(\sqrt{G})^{-1}y$, which is indeed $y$. Therefore, for any $y\in \mathbb{R}^m$, there exists an $x\in \mathbb{R}^m$ such that $f(x)=y$. As such, we have shown
                $f$ to be bijective.\newline
                We now show that $f$ satisfies the relationship $v\cdot_{G} w = f(v)\cdot f(w)$ for any $v,w \in \mathbb{R}^m$. By definition, $f(v)\cdot f(w)=(\sqrt{G}v)^{T}I\sqrt{G}w$. Using rules of transposition,
                this becomes $v^T\sqrt{G}I\sqrt{G}w$, which is indeed $v^TGw=v\cdot_{G} w$. Thus, we are done.
              \end{proof}
    \end{enumerate}
\end{enumerate}
\end{document}